\section{Lean KL Bridge}\label{app:lean}

The main text uses the expected-Kullback--Leibler form of closure deficit in \Cref{prop:cd-kl}. As a small formal appendix, the companion repository at \url{https://github.com/ioannist/six-birds-randomness} includes a machine-checked Lean development that bridges probability-level KL divergence to the expected log-likelihood ratio form used in the finite arguments \citep{demoura2015leantheoremprover,mathlib2020leanmathematicallibrary}. The formal artifact is intentionally modest. It does not attempt a full mechanization of conditional mutual information or of the complete closure-deficit theory. Instead it certifies the KL step on which the expected-divergence reading rests.

For probability measures $\mu$ and $\nu$ with $\mu \ll \nu$, the formal development proves the identity
\[
(\mathrm{klDiv}(\mu,\nu))^{\mathrm{toReal}}
=
\int \mathrm{llr}_{\mu,\nu}\, d\mu,
\]
that is, the real-valued Kullback--Leibler divergence equals the integral of the log-likelihood ratio under $\mu$. From that bridge it derives the nonnegativity statement
\[
0 \le \int \mathrm{llr}_{\mu,\nu}\, d\mu.
\]
The development also records the zero-divergence characterization for probability measures, namely that KL divergence vanishes exactly when the two measures are equal.

Conceptually, this appendix matters because the paper's finite-state formula
\[
\mathsf{CD}_\tau(\Pi)
=
\sum_x \pi(x)\,
D_{\mathrm{KL}}\!\bigl(p_x^{(\tau)} \,\|\, \bar p_{\Pi(x)}^{(\tau)}\bigr)
\]
is an instance of this more general expected-divergence pattern. The Lean artifact therefore does not formalize the entire manuscript, but it does machine-check the bridge between divergence and expected log-likelihood ratio that underwrites the KL reading of hidden predictive residue.

The tracked Lean files in the repository are the smoke module and the KL bridge module. Their role in the present paper is evidentiary rather than expository: they certify a core algebraic step, while the manuscript itself remains focused on the conceptual and empirical account of randomness as predictive non-closure.
